% ==============================================================================
% CHAPTER 6: EXPERIMENTS AND EVALUATION
% ==============================================================================
\chapter{Experiments and Evaluation}
\label{chap:experiments}

\section{Experimental Setup and Testbed Environment}
\label{sec:experimental_setup}

To rigorously evaluate the ADK framework and the Code-First paradigm, empirical experiments were conducted on an x86-64 workstation featuring an 8-core CPU, 32 GB RAM, and Python 3.12 under Windows and Linux virtualization. Inference orchestration was evaluated using state-of-the-art foundation models (Claude 3.5 Sonnet and Gemini 1.5/2.0 Pro) accessed via authenticated REST endpoints.

\section{Case Study: Microservice Security Analysis Framework}
\label{sec:case_study}

As the benchmark validation scenario, the system was tasked with synthesizing a complete end-to-end software project and corresponding thesis documentation for the topic: \emph{"Security Analysis and Vulnerability Assessment in Microservice Architectures"}.

The multi-agent swarm executed the complete workflow:
\begin{enumerate}
    \item Decomposed functional security requirements into MoSCoW deliverables (token auditing, network policy verification, anomaly detection).
    \item Identified and verified 18 SOTA scientific papers published between 2023 and 2026.
    \item Synthesized Python scanning modules, AST inspection routines, and a comprehensive \texttt{pytest} suite.
    \item Ran load and latency profiling under concurrent workloads ($1$ to $100$ simulated threads).
    \item Successfully compiled the full thesis package in LaTeX and Typst.
\end{enumerate}

\section{Quality Gates Evaluation Results}
\label{sec:gate_results}

The complete artifact collection was audited by the \texttt{MasterVerificationSuite}. Table \ref{tab:verification_results} details the evaluation metrics and gate statuses.

\begin{table}[H]
\centering
\small
\caption{Comprehensive evaluation results across the 7 Quality Gates of ADK}
\label{tab:verification_results}
\begin{tabularx}{\textwidth}{l l c c X}
\toprule
\textbf{Quality Gate} & \textbf{Evaluation Metric} & \textbf{Threshold} & \textbf{Observed Value} & \textbf{Status} \\
\midrule
\textbf{1. CodeVerification} & AST syntax \& Unit tests pass rate & $100\%$ passed & $28/28$ passed ($100\%$) & \textbf{PASS} \\
\textbf{2. MutationTesting} & Mutation Score ($MS$) & $\ge 60.0\%$ & $74.2\%$ & \textbf{PASS} \\
\textbf{3. CitationVerification} & SOTA recency \& DOI validity & $\text{year} \ge 2023$ & $100\%$ valid SOTA & \textbf{PASS} \\
\textbf{4. EnglishNaming} & ASCII English file \& code naming & $100\%$ compliant & $100\%$ compliant & \textbf{PASS} \\
\textbf{5. CrossConsistency} & AST symbol references in text & $0$ hallucinations & $0$ unmatched symbols & \textbf{PASS} \\
\textbf{6. AcademicStyle} & Generative AI cliché count & $0$ occurrences & $0$ clichés detected & \textbf{PASS} \\
\textbf{7. StylometryAudit} & Type-Token Ratio ($TTR$) & $\ge 0.35$ & $0.412$ & \textbf{PASS} \\
\bottomrule
\end{tabularx}
\end{table}

All seven quality gates passed successfully on the final generation pass, confirming that the synthesized software and documentation meet rigorous production and academic standards.

\section{Mutation Analysis and Fault Detection}
\label{sec:mutation_evaluation}

A notable finding emerged from the comparison between line coverage and mutation testing. The initial test suite authored by the \texttt{DeveloperAgent} achieved a high line coverage of $91.4\%$. However, the initial \texttt{MutationTestingGate} run revealed a mutation score of only $54.1\%$ (failing the $60\%$ threshold) due to missing boundary assertions in network timeout handling.

Upon receiving the mutation audit report, the agent's self-repair loop generated targeted assertion tests covering edge conditions, killing 19 previously surviving mutants and raising the final mutation score to $74.2\%$. This demonstrates the vital necessity of mutation testing in agentic software development.

\section{Empirical Latency and Throughput Benchmarking}
\label{sec:latency_benchmarks}

The \texttt{ExperimenterAgent} executed automated load tests measuring the processing latency of the vulnerability scanning pipeline under increasing concurrency. The observed metrics are summarized in Table \ref{tab:sla_benchmarks}.

\begin{table}[H]
\centering
\small
\caption{Service latency percentiles ($p95$, $p99$) under varying concurrent workloads}
\label{tab:sla_benchmarks}
\begin{tabular}{ccccc}
\toprule
\textbf{Concurrency (Threads)} & \textbf{Mean [ms]} & \textbf{Median [ms]} & \textbf{95th Percentile [ms]} & \textbf{99th Percentile [ms]} \\
\midrule
1 & 12.4 & 11.8 & 16.2 & 21.0 \\
10 & 18.7 & 17.2 & 28.5 & 39.4 \\
50 & 44.1 & 41.5 & 68.9 & 92.3 \\
100 & 91.2 & 86.0 & 142.7 & 185.1 \\
\bottomrule
\end{tabular}
\end{table}

These measurements were captured directly by the runtime telemetry engine and programmatically injected into the thesis source, guaranteeing absolute consistency between benchmark telemetry and documented results.

